\section{Kernel-Induced Functional Representations}\label{sec:functional_representations}

Beyond regression, the covariance kernel induces functional representations exploited later in this thesis. \Cref{sec:reduced_rank} replaces the GP by a finite-dimensional projection to alleviate the cubic cost of kriging, and \cref{sec:RKHS} interprets kriging as minimum-norm interpolation in a reproducing kernel Hilbert space.

\subsection{Reduced-Rank Approximation}\label{sec:reduced_rank}

The dominant cost of kriging is the inversion of the $\tilde{n}\times\tilde{n}$ Gram matrix $G$, which scales as $\mathcal{O}(\tilde{n}^3)$ and becomes prohibitive for large training sets. Reduced-rank methods approximate the kernel by a finite sum $k(x,x')\approx\sum_{j=1}^r s_j\:\phi_j(x)\:\phi_j(x')$ over $r\ll\tilde{n}$ fixed basis functions, which turns the $\tilde{n}\times\tilde{n}$ inversion into an $r\times r$ one. We follow the Hilbert-space method, tailored to stationary kernels, whose basis does not depend on the kernel hyperparameters.

Assume $k$ to be stationary and radial on $\mathcal{U}=\mathbb{R}^{d_x}$, \textit{i.e.} $k(x,x')=k(\|x-x'\|)$. By Bochner's theorem, $k$ is then the inverse Fourier transform of a non-negative spectral density $S$, a duality known as the Wiener-Khinchin theorem:
\begin{equation}
	k(r)=\frac{1}{(2\pi)^{d_x}}\int_{\mathbb{R}^{d_x}}S(\omega)\:e^{i\omega^\top r}\:d\omega,\qquad S(\omega)=\int_{\mathbb{R}^{d_x}}k(r)\:e^{-i\omega^\top r}\:dr
\end{equation}
For the squared exponential kernel \eqref{eq:SE_kernel}, the spectral density is itself a Gaussian:
\begin{equation}
	S(\omega)=\sigma_f^2\:(2\pi)^{d_x/2}\left(\prod_{i=1}^{d_x}l_i\right)\exp\left(-\frac{1}{2}\omega^\top\textup{diag}(l)^2\:\omega\right)
\end{equation}

Restricting the field to a bounded domain $\Omega\subset\mathbb{R}^{d_x}$, let $(\lambda_j,\phi_j)_{j\geq1}$ be the eigenpairs of the negative Laplacian on $\Omega$ with Dirichlet boundary conditions:
\begin{equation}
	-\nabla^2\phi_j=\lambda_j\:\phi_j\:\text{ on }\Omega,\qquad\phi_j=0\:\text{ on }\partial\Omega
\end{equation}
which form an orthonormal basis of $L^2(\Omega)$. The stationary covariance operator $\varphi\mapsto\int k(\cdot,x')\:\varphi(x')\:dx'$ is a function of the Laplacian whose transfer function is the spectral density $S$; replacing the frequency $\|\omega\|^2$ by the Laplacian eigenvalues $\lambda_j$ yields the approximation:
\begin{equation}
	k(x,x')\approx\sum_{j=1}^rS\left(\sqrt{\lambda_j}\right)\phi_j(x)\:\phi_j(x')\label{eq:reduced_rank_kernel}
\end{equation}
On a box domain $\Omega=\prod_{i=1}^{d_x}[-L_i,L_i]$, the eigenpairs are explicit, indexed by $j=(j_1,\dots,j_{d_x})\in(\mathbb{N}^*)^{d_x}$:
\begin{equation}
	\phi_j(x)=\prod_{i=1}^{d_x}\frac{1}{\sqrt{L_i}}\sin\left(\frac{\pi\:j_i\:(x_i+L_i)}{2L_i}\right),\qquad\lambda_j=\sum_{i=1}^{d_x}\left(\frac{\pi\:j_i}{2L_i}\right)^2
\end{equation}

Collecting the basis functions in $\Phi(x)\triangleq\begin{bmatrix}\phi_1(x)&\cdots&\phi_r(x)\end{bmatrix}^\top$ and the spectral weights in the diagonal matrix $\Lambda\triangleq\textup{diag}\left(S(\sqrt{\lambda_1}),\dots,S(\sqrt{\lambda_r})\right)$, \eqref{eq:reduced_rank_kernel} reads $k(x,x')\approx\Phi(x)^\top\Lambda\:\Phi(x')$. The regularised Gram matrix thus becomes $G\approx\Phi_{\tilde{X}}\:\Lambda\:\Phi_{\tilde{X}}^\top+\tilde{R}\:I_{\tilde{n}}$, where $\Phi_{\tilde{X}}\triangleq\begin{bmatrix}\Phi(\tilde{x}_1)&\cdots&\Phi(\tilde{x}_{\tilde{n}})\end{bmatrix}^\top\in\mathcal{M}_{\tilde{n},r}(\mathbb{R})$. As this is a rank-$r$ perturbation of a multiple of the identity, the Woodbury inversion formula (\Cref{th:Woodbury}) reduces the $\tilde{n}\times\tilde{n}$ inversion to an $r\times r$ one, and the simple kriging estimator \eqref{eq:simple_kriging_equations} becomes, in the centred case:
\begin{equation}
	z_s^*\approx\Phi(x)^\top\left(\Phi_{\tilde{X}}^\top\Phi_{\tilde{X}}+\tilde{R}\:\Lambda^{-1}\right)^{-1}\Phi_{\tilde{X}}^\top\tilde{Z}
\end{equation}
bringing the overall cost down from $\mathcal{O}(\tilde{n}^3)$ to $\mathcal{O}(r^2\tilde{n})$.

\begin{remark}[Hyperparameter-independent basis]
	A decisive advantage of \eqref{eq:reduced_rank_kernel} is that the basis functions $\phi_j$, being eigenfunctions of the Laplacian, do not depend on the kernel hyperparameters $\theta$: only the diagonal weights $S(\sqrt{\lambda_j})$ do. Evidence maximisation (\Cref{th:GP_MLE}) thus recomputes only the diagonal $\Lambda$, the costly product $\Phi_{\tilde{X}}^\top\Phi_{\tilde{X}}$ being formed once and for all. The approximation becomes exact as the rank $r$ and the domain $\Omega$ grow to infinity; for the squared exponential kernel, the truncation error even decreases at a rate independent of the input dimension $d_x$.
\end{remark}

\begin{remark}[Weight-space view]
	Equation \eqref{eq:reduced_rank_kernel} amounts to the Bayesian linear model $f(x)=\Phi(x)^\top w$ with weight prior $w\sim\mathcal{N}(0,\Lambda)$. This weight-space formulation is the starting point for the random Fourier features and kernel recursive least-squares methods used for the SLAM application of this thesis.
\end{remark}

\subsection{Reproducing Kernel Hilbert Space Formulation}\label{sec:RKHS}

The covariance kernel also induces a deterministic, function-analytic counterpart to GP regression. Let $k:\mathcal{U}\times\mathcal{U}\to\mathbb{R}$ be a scalar SPSD kernel on an arbitrary index set $\mathcal{U}$.

\begin{definition}[Reproducing kernel Hilbert space]\label{th:RKHS}
	The reproducing kernel Hilbert space (RKHS) $\mathcal{H}_k$ associated with $k$ is the completion of the span of $\{k(\cdot,x)\:|\:x\in\mathcal{U}\}$ for the inner product defined on the generators by:
	\begin{equation}
		\langle k(\cdot,x)\:|\:k(\cdot,x')\rangle_{\mathcal{H}_k}\triangleq k(x,x')
	\end{equation}
	It is characterised by the reproducing property: for all $\psi\in\mathcal{H}_k$ and $x\in\mathcal{U}$,
	\begin{equation}
		\psi(x)=\langle\psi\:|\:k(\cdot,x)\rangle_{\mathcal{H}_k}\label{eq:reproducing_property}
	\end{equation}
\end{definition}

The reproducing property \eqref{eq:reproducing_property} states that point evaluation is a continuous linear functional on $\mathcal{H}_k$, which endows the space with a regularity entirely dictated by $k$. This space provides a variational reading of kriging.

\begin{proposition}[Kriging as minimum-norm interpolation]\label{th:RKHS_kriging}
	In the centred case, the simple kriging estimator \eqref{eq:simple_kriging_equations} is the solution of the penalised least-squares problem in $\mathcal{H}_k$:
	\begin{equation}
		z_s^*=\underset{\psi\in\mathcal{H}_k}{\textup{argmin}}\left(\frac{1}{\tilde{R}}\sum_{i=1}^{\tilde{n}}\left(\psi(\tilde{x}_i)-\tilde{z}_i\right)^2+\|\psi\|_{\mathcal{H}_k}^2\right)\label{eq:RKHS_regularisation}
	\end{equation}
	which, in the noise-free limit $\tilde{R}\to0$, reduces to the minimum-norm interpolant $\textup{argmin}\{\|\psi\|_{\mathcal{H}_k}\:|\:\psi(\tilde{x}_i)=\tilde{z}_i\}$.
\end{proposition}

\begin{proof}
	By the representer theorem, a minimiser of \eqref{eq:RKHS_regularisation} lies in the finite-dimensional span of $\{k(\cdot,\tilde{x}_i)\}_{i=1}^{\tilde{n}}$, \textit{i.e.} $\psi=\sum_i\alpha_i\:k(\cdot,\tilde{x}_i)$. The reproducing property gives $\psi(\tilde{x}_j)=\sum_i\alpha_i\:k(\tilde{x}_j,\tilde{x}_i)$ and $\|\psi\|_{\mathcal{H}_k}^2=\alpha^\top k(\tilde{X},\tilde{X})\:\alpha$, so \eqref{eq:RKHS_regularisation} becomes the quadratic problem $\textup{argmin}_\alpha\:\tilde{R}^{-1}\|k(\tilde{X},\tilde{X})\alpha-\tilde{Z}\|^2+\alpha^\top k(\tilde{X},\tilde{X})\alpha$. Cancelling its gradient gives $\alpha=G^{-1}\tilde{Z}$, hence $\psi(x)=k(x,\tilde{X})\:G^{-1}\tilde{Z}$, which is \eqref{eq:simple_kriging_equations} for $m=0$.
\end{proof}

\begin{remark}[GP and RKHS correspondence]
	The two viewpoints are complementary: the GP model gives a probabilistic reading of kriging (posterior mean and variance under a Gaussian prior), while the RKHS gives a deterministic one (minimum-norm interpolant in a function space). Writing the Mercer decomposition of the kernel over its eigenpairs $(\beta_j,\phi_j)_{j\geq1}$ as $k(x,x')=\sum_{j\geq1}\beta_j\:\phi_j(x)\:\phi_j(x')$, a function $\psi=\sum_{j\geq1}c_j\:\phi_j$ belongs to $\mathcal{H}_k$ with squared norm $\|\psi\|_{\mathcal{H}_k}^2=\sum_{j\geq1}c_j^2/\beta_j$; the norm thus penalises most the components carried by low-eigenvalue, rapidly oscillating eigenfunctions, so the decay of the eigenvalues $\beta_j$ sets the smoothness enforced by $k$. Choosing a kernel therefore amounts to choosing a regularity class for the unknown field: a smooth kernel such as the SE kernel yields infinitely differentiable interpolants, whereas a rougher Matérn kernel admits functions of limited smoothness.
\end{remark}
